\subsection*{Recovering a phrase in Song verse}
We next asked whether a writer that distinguishes two plausible sources also recovers the expression that the poet wrote. Four continuations of the Tang models saw identical batches of line pairs. All four learned the ordered probability of two adjacent missing glyphs. Two also learned to distinguish a poet's authentic source--ending pairs from exchanged pairs. We tested on 3,854 Song line pairs from 192 poet attributions absent from the Tang corpus (Fig.~\ref{fig:verse-contrast}).

When given both authentic endings, the contrastively trained image writer recovered the original assignment on 3,020 of 3,781 scorable comparisons, compared with 2,970 for its matched control. When asked to write the two missing cells, it recovered eight recorded phrases, compared with 12 for the control and four for the frozen parent. All eight successes were shared with the control. The paired difference from the control was $-0.10$ percentage points (author-bootstrap 95\% interval $-0.22$ to $0.00$). Exchanging the pictured first line reduced the eight exact repairs to two. The matched contrastive byte writer completed 17 repairs. None of the six neural writers recovered a complete line starting with all target cells blank.

Only 38 pages required an output glyph absent from the bank. On 555 pages the contrastive image writer placed the correct first glyph among its eight candidates. Even when supplied with the true first glyph, it ranked the recorded second glyph first on only 151 pages. Both requirements held together on 28 pages, counting ties in favour of the reference. The results locate a limitation in the learned conditional distribution. A wider search could admit additional first glyphs, but even exhaustive first-glyph search with this second-cell rule could recover at most 151 recorded pairs. Exact restoration credits one historical reading and cannot determine the literary merits of alternatives.
